\documentclass[preprint,review,12pt]{article}

% Packages for JMVA-style formatting
\usepackage[utf8]{inputenc}
\usepackage[T1]{fontenc}
\usepackage{amsmath,amssymb,amsthm}
\usepackage{mathtools}
\usepackage{enumitem}
\usepackage{hyperref}
\usepackage{graphicx}
\usepackage{booktabs}
\usepackage{array}
\usepackage[margin=2.5cm]{geometry}
\usepackage{natbib}
\usepackage{lineno}
\usepackage{authblk}
\usepackage{algorithm}
\usepackage{algorithmic}

\linenumbers

\setcitestyle{authoryear,round}

% Theorem environments
\newtheorem{theorem}{Theorem}[section]
\newtheorem{lemma}[theorem]{Lemma}
\newtheorem{proposition}[theorem]{Proposition}
\newtheorem{corollary}[theorem]{Corollary}
\newtheorem{definition}[theorem]{Definition}
\newtheorem{remark}[theorem]{Remark}
\newtheorem{example}[theorem]{Example}
\newtheorem{openproblem}[theorem]{Open Problem}

% Custom commands
\newcommand{\R}{\mathbb{R}}
\newcommand{\N}{\mathbb{N}}
\newcommand{\E}{\mathbb{E}}
\newcommand{\Var}{\mathrm{Var}}
\newcommand{\KL}{\mathrm{KL}}

\begin{document}
	
	\title{Minimum Description Length Estimation for Stochastic Processes: \\
		A Unified Framework for Statistical Learning and Data Compression}
	
	\author[1,2]{Jocelyn Nembé}
	\affil[1]{L.I.A.G.E, Institut National des Sciences de Gestion, BP 190, Libreville, Gabon}
	\affil[2]{Modeling and Calculus Lab, ROOTS-INSIGHTS, Libreville, Gabon}
	\affil[ ]{Email: jnembe@hotmail.com, jnembe@root-insights.com}
	
	\date{December 2025}
	
	\maketitle
	
	\begin{abstract}
		We develop a unified framework for consistent estimation across a broad class 
		of stochastic processes using the minimum description length (MDL) principle. 
		Our primary contribution is methodological: we demonstrate that a single 
		estimation procedure achieves consistent inference for fundamentally different 
		stochastic structures, including density estimation, Poisson processes, point 
		processes with multiplicative intensities, continuous-time Markov chains, and 
		diffusion processes. For each process class, we derive explicit relationships 
		between the Hellinger distance of probability measures and the corresponding 
		functional parameters, revealing a common mathematical structure that enables 
		unified analysis. The MDL estimator satisfies an oracle inequality of the form 
		$H^2(P, \hat{P}_n) \leq \inf_{Q \in \Gamma_n}(H^2(P,Q) + C_n(Q)/n) + O((\log n)/n)$,
		automatically balancing approximation accuracy and model complexity. We develop 
		practical applications to universal data compression and automatic model 
		selection, providing explicit algorithms that handle heterogeneous data sources 
		without prior knowledge of the underlying stochastic structure.
	\end{abstract}
	
	\noindent\textbf{Keywords:} Minimum description length; Nonparametric estimation; 
	Hellinger distance; Counting processes; Diffusion processes; Point processes; 
	Model selection; Data compression; Universal coding
	
	\noindent\textbf{MSC 2020:} 60G55; 62G05; 62G20; 62M05; 94A29
	
	\vspace{1em}
	\noindent\textbf{Highlights:}
	\begin{itemize}[leftmargin=2em]
		\item Unified MDL estimator applicable across five classes of stochastic processes
		\item Explicit Hellinger distance formulas linking measure-level and parameter-level distances
		\item Oracle inequality demonstrating automatic bias-complexity tradeoff
		\item Practical algorithms for universal compression and automatic model selection
		\item Theoretical foundation for handling heterogeneous data with a single procedure
	\end{itemize}
	
	%==============================================================================
	\section{Introduction}
	\label{sec:intro}
	%==============================================================================
	
	%------------------------------------------------------------------------------
	\subsection{Motivation and Objectives}
	%------------------------------------------------------------------------------
	
	Modern data sources exhibit remarkable heterogeneity. A single application 
	may need to process continuous signals, event sequences, state transitions, 
	and diffusion trajectories---each traditionally requiring distinct statistical 
	methodologies. Video compression, for instance, must handle both the temporal 
	dynamics of scene changes (naturally modeled as point processes) and the 
	spatial correlations within frames (often modeled through Markovian or 
	diffusion-type structures). Similarly, in machine learning pipelines, 
	practitioners frequently encounter mixed data types that do not fit neatly 
	into a single probabilistic framework.
	
	This fragmentation poses practical challenges. Developing separate estimation 
	procedures for each data type leads to:
	\begin{itemize}[leftmargin=2em]
		\item duplicated algorithmic efforts across application domains,
		\item difficulty in building unified software systems that handle 
		heterogeneous data streams,
		\item lack of principled methods for model selection when the underlying 
		stochastic structure is unknown or evolves over time.
	\end{itemize}
	
	The present work addresses these challenges by developing a \emph{unified 
		estimation framework} based on the Minimum Description Length (MDL) principle. 
	Our primary objective is not to establish new convergence rates---which are 
	well-understood for individual model classes---but rather to demonstrate that 
	a \emph{single estimation principle} achieves consistent inference across 
	fundamentally different stochastic structures, with explicit and comparable 
	performance guarantees.
	
	%------------------------------------------------------------------------------
	\subsection{The Unification Problem}
	%------------------------------------------------------------------------------
	
	Consider the following estimation problems, each central to different 
	application domains:
	
	\begin{enumerate}[label=(\roman*)]
	\item \textbf{Density estimation:} Given i.i.d.\ samples $X_1, \ldots, X_n$ 
	from an unknown density $p$, estimate $p$.
	
	\item \textbf{Poisson process estimation:} Given event times 
	$(T_1, \ldots, T_N)$ from a Poisson process with unknown mean measure 
	$\mu$, estimate $\mu$.
	
	\item \textbf{Point process intensity estimation:} Given observations of 
	a counting process $N(t)$ with multiplicative intensity 
	$\lambda(t) = \alpha(t) Y(t)$, estimate the baseline intensity $\alpha$.
	
	\item \textbf{Markov transition estimation:} Given sample paths of a 
	continuous-time Markov chain, estimate the transition intensities 
	$\alpha^{hj}(t)$.
	
	\item \textbf{Diffusion drift estimation:} Given a trajectory of the 
	process $dX_t = \theta(X_t) dt + \sigma dW_t$, estimate the drift 
	function $\theta$.
\end{enumerate}

Each problem has been extensively studied using tailored techniques: kernel 
methods and sieves for densities \citep{Devroye2001}, martingale estimating 
equations for point processes \citep{Andersen1993}, maximum likelihood and 
Girsanov transformations for diffusions \citep{Kutoyants2004}. However, these 
approaches share a common mathematical structure that has not been fully 
exploited: in each case, estimation reduces to optimizing a likelihood-type 
criterion over a space of candidate models, and consistency is governed by 
the metric geometry of this space.

The key insight underlying our unified framework is that the Hellinger 
distance provides a \emph{universal metric} for measuring estimation accuracy 
across all these settings. While the functional parameters differ (densities, 
intensities, drift functions), the probability measures they induce can all 
be compared through Hellinger distance, and this comparison admits explicit 
formulas linking the measure-level distance to the parameter-level distance.

%------------------------------------------------------------------------------
\subsection{Contributions}
%------------------------------------------------------------------------------

This paper makes three main contributions:

\paragraph{1. A unified MDL estimation principle.}
We formulate a single minimum complexity estimator that applies uniformly to 
densities, Poisson processes, general point processes, Markov chains, and 
diffusions. The estimator takes the form
\begin{equation}\label{eq:unified_estimator}
	\hat{P}_n = \arg\min_{Q \in \Gamma_n} \left( -\log \frac{dQ}{dP_0}
	(\text{data}) + C_n(Q) \right),
\end{equation}
where $\Gamma_n$ is a countable model class, $P_0$ is a reference measure, 
and $C_n(Q)$ is a complexity penalty satisfying the Kraft inequality. The 
same algorithmic template---enumerate models, compute penalized 
log-likelihoods, select the minimizer---applies across all process types.

\paragraph{2. Explicit Hellinger distance relationships.}
For each process class, we derive the precise relationship between the 
Hellinger distance of probability measures and the natural distance on 
functional parameters. These relationships take the form:
\begin{equation}\label{eq:hellinger_template}
	H^2(P_\theta, P_\gamma) = 1 - \mathbb{E}_\zeta\left[ 
	\exp\left( -c \int d^2(\theta, \gamma) \, d\nu \right) \right],
\end{equation}
where $\zeta$ is an intermediate parameter, $c$ is a model-specific constant, 
$d(\theta, \gamma)$ is the $L^2$-type distance between parameters, and $\nu$ 
is an appropriate measure. Table~\ref{tab:hellinger_summary} summarizes these 
relationships. Crucially, all formulas share similar analytic properties, 
enabling a unified convergence analysis.

\paragraph{3. Unified oracle inequality.}
We establish that the MDL estimator satisfies, across all process classes,
\begin{equation}\label{eq:oracle_inequality}
	H^2(P, \hat{P}_n) \leq \inf_{Q \in \Gamma_n} \left( H^2(P, Q) 
	+ \frac{C_n(Q)}{n} \right) + O\left( \frac{\log n}{n} \right),
\end{equation}
where the approximation term $H^2(P, Q)$ measures model misspecification 
and the complexity term $C_n(Q)/n$ penalizes model richness. This oracle 
inequality demonstrates that MDL estimation automatically balances the 
bias-variance tradeoff, achieving near-optimal performance without requiring 
prior knowledge of the true model's complexity.

%------------------------------------------------------------------------------
\subsection{Significance for Applications}
%------------------------------------------------------------------------------

The unified framework has direct implications for two application domains:

\paragraph{Universal data compression.}
The MDL principle originates from the observation that optimal statistical 
inference and optimal data compression are fundamentally equivalent 
\citep{Rissanen1978, Cover2006}. An estimator achieving the bound 
\eqref{eq:oracle_inequality} automatically yields a two-stage code with 
redundancy
\[
R_n = \frac{1}{n}\left( -\log \hat{P}_n(\text{data}) + \log P(\text{data})
\right) \leq \inf_{Q \in \Gamma_n} \left( D(P \| Q) + \frac{C_n(Q)}{n}
\right),
\]
where $D(P \| Q)$ is the Kullback-Leibler divergence. Our results imply that 
a \emph{single compression algorithm} based on MDL can achieve near-optimal 
compression rates for data generated by any of the five process classes 
considered, without prior knowledge of which class generated the data.

\paragraph{Automatic model selection.}
In many applications, the practitioner does not know \emph{a priori} whether 
data should be modeled as i.i.d., as a point process, or as a diffusion. 
The MDL framework provides a principled solution: include models from 
multiple classes in $\Gamma_n$, assign appropriate complexity penalties, 
and let the data select the best representation. The oracle inequality 
\eqref{eq:oracle_inequality} guarantees that this procedure performs nearly 
as well as if the correct model class had been known in advance.

%------------------------------------------------------------------------------
\subsection{Relation to Prior Work}
%------------------------------------------------------------------------------

The theoretical foundations of MDL estimation were laid by 
\citet{Rissanen1978, Rissanen1983, Rissanen1984} and developed into a 
comprehensive framework by \citet{Barron1991}, \citet{Barron1998}, and 
\citet{Grunwald2007}. Oracle inequalities of the type \eqref{eq:oracle_inequality} 
for density estimation appear in \citet{Barron1999} and \citet{Yang1999}, 
with connections to model selection explored by \citet{Birge1998, Birge2006} 
and \citet{Hansen2001}.

For specific process classes, sieve estimation and penalized likelihood 
methods have been developed by \citet{Karr1987} and \citet{McKeague1986} 
for point processes, by \citet{Andersen1993} for counting processes, and 
by \citet{Kutoyants2004} and \citet{PrakasaRao1999a} for diffusions. The 
role of Hellinger distance in nonparametric estimation was clarified by 
\citet{LeCam1973, LeCam1986} and \citet{Birge1983}.

\textbf{Our contribution relative to this literature} is not to improve 
convergence rates for any individual model class---the rates we obtain are 
comparable to known results---but rather to provide:
\begin{enumerate}[label=(\alph*)]
\item a \emph{single, explicit estimator} that works across all classes 
without modification,
\item \emph{explicit formulas} linking Hellinger distances at the measure 
level to parameter distances, enabling direct comparison of estimation 
difficulty across model types,
\item a \emph{unified proof technique} that highlights the common 
mathematical structure underlying consistency in heterogeneous settings.
\end{enumerate}

To our knowledge, no prior work has systematically developed MDL estimation 
across densities, Poisson processes, general point processes, Markov chains, 
and diffusions within a single framework with explicit, comparable bounds.

%------------------------------------------------------------------------------
\subsection{Paper Organization}
%------------------------------------------------------------------------------

Section~\ref{sec:setting} introduces the statistical setting, defines the 
Hellinger distance and affinity, and establishes their key properties for 
product measures. Section~\ref{sec:convergence} presents the main convergence 
theorem for MDL estimators in abstract form. Section~\ref{sec:applications} 
instantiates this theorem for each of the five process classes, deriving 
the explicit Hellinger relationships and stating the corresponding 
convergence results, with applications to compression and model selection. 
Section~\ref{sec:proofs} contains all proofs. Section~\ref{sec:discussion} 
discusses limitations, extensions, and open problems.

%------------------------------------------------------------------------------
% Summary table of Hellinger relationships
%------------------------------------------------------------------------------

\begin{table}[t]
\centering
\caption{Summary of Hellinger distance relationships across process classes. 
	In each case, $H^2(P_\theta, P_\gamma)$ denotes the squared Hellinger distance 
	between probability measures, and the right column shows the explicit formula 
	in terms of functional parameters.}
\label{tab:hellinger_summary}
\begin{tabular}{@{}lll@{}}
	\toprule
	\textbf{Process Class} & \textbf{Parameter} & \textbf{Hellinger Relationship} \\
	\midrule
	I.i.d.\ density & density $p$ & 
	$H^2(P^n, Q^n) = 1 - (1 - H^2(P,Q))^n$ \\[6pt]
	Poisson process & intensity $\alpha$ & 
	$H^2(P_\alpha, P_\beta) = 1 - e^{-H^2(\alpha, \beta)}$ \\[6pt]
	Point process & intensity $\alpha$ & 
	$H^2(P_\alpha, P_\beta) = 1 - \mathbb{E}\left[e^{-\frac{1}{2}\int 
		(\sqrt{\alpha} - \sqrt{\beta})^2 Y \, ds}\right]$ \\[6pt]
	Markov chain & transitions $\alpha^{hj}$ & 
	$H^2(P_\alpha, P_\gamma) = 1 - \mathbb{E}\left[e^{-\frac{1}{2}\sum_{h,j}\int 
		(\sqrt{\alpha^{hj}} - \sqrt{\gamma^{hj}})^2 Y_h \, ds}\right]$ \\[6pt]
	Diffusion & drift $\theta$ & 
	$H^2(P_\theta, P_\gamma) = 1 - \mathbb{E}\left[e^{-\frac{1}{8\sigma^2}
		\int (\theta - \gamma)^2 dt}\right]$ \\
	\bottomrule
\end{tabular}
\end{table}

%==============================================================================
\section{The Statistical Setting}
\label{sec:setting}
%==============================================================================

%------------------------------------------------------------------------------
\subsection{Hellinger Distance and Affinity}
%------------------------------------------------------------------------------

For two probability measures $P$ and $Q$ on a measurable space $(\mathcal{X}, \mathcal{A})$, 
both dominated by a $\sigma$-finite measure $\mu$, the \emph{Hellinger distance} 
$H(P, Q)$ is defined as \citep{LeCam1986}:
\begin{equation}\label{eq:hellinger_def}
H^2(P, Q) = \frac{1}{2} \int_{\mathcal{X}} \left( \sqrt{\frac{dP}{d\mu}} 
- \sqrt{\frac{dQ}{d\mu}} \right)^2 d\mu.
\end{equation}

The \emph{Hellinger affinity} between $P$ and $Q$ is defined as \citep{Kakutani1948}:
\begin{equation}\label{eq:affinity_def}
\rho(P, Q) = \int_{\mathcal{X}} \sqrt{\frac{dP}{d\mu} \frac{dQ}{d\mu}} \, d\mu.
\end{equation}

These quantities are related by:
\begin{equation}\label{eq:affinity_distance}
H^2(P, Q) = 1 - \rho(P, Q).
\end{equation}

\begin{remark}[Convention]
Throughout this paper, we use the convention that $H^2(P,Q) \in [0,1]$ with 
$H^2(P,Q) = 0$ if and only if $P = Q$, and $H^2(P,Q) = 1$ if and only if 
$P$ and $Q$ are mutually singular. This corresponds to the ``squared Hellinger 
distance'' in some references.
\end{remark}

%------------------------------------------------------------------------------
\subsection{Properties of Hellinger Distance}
%------------------------------------------------------------------------------

The Hellinger distance enjoys several important properties:

\begin{lemma}[Basic Properties]\label{lem:basic_properties}
The Hellinger distance satisfies:
\begin{enumerate}[label=(\roman*)]
	\item \textbf{Metric:} $H$ is a metric on the space of probability measures.
	\item \textbf{Bounds:} $0 \leq H^2(P, Q) \leq 1$.
	\item \textbf{Relation to total variation:} $H^2(P,Q) \leq \|P - Q\|_{TV} \leq \sqrt{2} H(P,Q)$.
	\item \textbf{Relation to KL divergence:} $2H^2(P, Q) \leq D(P \| Q)$.
\end{enumerate}
\end{lemma}

%------------------------------------------------------------------------------
\subsection{Hellinger Distance for Product Measures}
%------------------------------------------------------------------------------

A crucial property for our analysis is the behavior of Hellinger distance 
under products.

\begin{lemma}[Product Measures]\label{lem:product_measures}
Let $P^n = P \otimes \cdots \otimes P$ and $Q^n = Q \otimes \cdots \otimes Q$ 
be $n$-fold product measures. Then:
\begin{equation}\label{eq:product_affinity}
	\rho(P^n, Q^n) = [\rho(P, Q)]^n = [1 - H^2(P, Q)]^n,
\end{equation}
and consequently:
\begin{equation}\label{eq:product_hellinger}
	H^2(P^n, Q^n) = 1 - [1 - H^2(P, Q)]^n.
\end{equation}
\end{lemma}

\begin{proof}
The affinity factorizes:
\begin{align*}
	\rho(P^n, Q^n) &= \int \prod_{i=1}^n \sqrt{\frac{dP}{d\mu}(x_i) 
		\frac{dQ}{d\mu}(x_i)} \, d\mu^n \\
	&= \prod_{i=1}^n \int \sqrt{\frac{dP}{d\mu} \frac{dQ}{d\mu}} \, d\mu 
	= [\rho(P,Q)]^n.
\end{align*}
The result follows from $H^2 = 1 - \rho$.
\end{proof}

\begin{corollary}[Approximation for Small Distances]\label{cor:small_distance}
For $H^2(P, Q) \ll 1$:
\begin{equation}\label{eq:small_approximation}
	H^2(P^n, Q^n) \approx n \cdot H^2(P, Q).
\end{equation}
\end{corollary}

This corollary explains why the complexity term in the oracle inequality 
\eqref{eq:oracle_inequality} is scaled by $1/n$: the information from $n$ 
observations accumulates linearly, while the complexity cost is fixed.

%------------------------------------------------------------------------------
\subsection{The MDL Estimation Principle}
%------------------------------------------------------------------------------

Let $\Gamma_n$ be a countable collection of candidate probability measures, 
and let $C_n: \Gamma_n \to [0, \infty)$ be a \emph{complexity function} 
satisfying the Kraft inequality:
\begin{equation}\label{eq:kraft}
\sum_{Q \in \Gamma_n} 2^{-C_n(Q)} \leq 1.
\end{equation}

\begin{remark}[Logarithmic base]\label{rem:log_base}
We adopt the convention that $\log$ denotes the natural logarithm throughout
the analysis, so that complexities $C_n(Q)$ are measured in nats and the Kraft
inequality \eqref{eq:kraft} reads $\sum_{Q} e^{-C_n(Q)} \leq 1$ in these units.
We retain the base-$2$ form in \eqref{eq:kraft} (and the bits unit in
Section~\ref{subsec:compression}) for its coding interpretation: a change of
base merely rescales $C_n(Q)$ by the constant factor $\log 2$ and is absorbed
into the constants of the oracle inequality. All exponential and Markov
inequalities below are stated with the natural base accordingly.
\end{remark}

The \emph{minimum description length (MDL) estimator} is defined as:
\begin{equation}\label{eq:mdl_estimator}
\hat{P}_n = \arg\min_{Q \in \Gamma_n} \left( -\log \frac{dQ}{dP_0}(X^n) 
+ C_n(Q) \right),
\end{equation}
where $P_0$ is a fixed reference measure dominating all $Q \in \Gamma_n$, 
and $X^n$ denotes the observed data.

The interpretation is:
\begin{itemize}
\item $-\log \frac{dQ}{dP_0}(X^n)$: code length for the data given model $Q$
\item $C_n(Q)$: code length for describing the model $Q$
\item The MDL estimator minimizes total description length
\end{itemize}

%==============================================================================
\section{Convergence of the MDL Estimator}
\label{sec:convergence}
%==============================================================================

%------------------------------------------------------------------------------
\subsection{Main Result}
%------------------------------------------------------------------------------

\begin{theorem}[Main Convergence Theorem]\label{thm:main}
Let $(\Omega^n, \mathcal{A}^n, P)$ be a probability space and $\Gamma_n$ a 
countable collection of probability measures with complexity function 
$C_n: \Gamma_n \to [0, \infty)$ satisfying the Kraft inequality 
\eqref{eq:kraft}. Let $P_0$ be a reference measure dominating all 
$Q \in \Gamma_n$.

The MDL estimator \eqref{eq:mdl_estimator} satisfies, for any $P$:
\begin{equation}\label{eq:main_result}
	H^2(P, \hat{P}_n) \leq \inf_{Q \in \Gamma_n} \left( H^2(P, Q) 
	+ \frac{C_n(Q)}{n} \right) + O\left( \frac{\log n}{n} \right)
\end{equation}
$P$-almost surely as $n \to \infty$.

More precisely, for any sequence $P_n \in \Gamma_n$:
\begin{equation}\label{eq:probability_bound}
	P\left( H^2(P, \hat{P}_n) > r_n \right) \leq 
	2 \exp\left( -\frac{n}{2} \left( r_n - H^2(P, P_n) - \frac{C_n(P_n)}{n} 
	\right)_+ \right) + o(1),
\end{equation}
where $(x)_+ = \max(x, 0)$.
\end{theorem}

The proof is given in Section~\ref{sec:proofs}.

\begin{remark}[On the complexity condition]\label{rem:half_kraft}
The proof of Theorem~\ref{thm:main} controls the union bound through the
quantity $\sum_{Q \in \Gamma_n} 2^{-C_n(Q)/2}$ (see Step~2 in
Section~\ref{sec:proofs}), and the argument as written requires the
\emph{summability at half exponent}
\begin{equation}\label{eq:half_kraft}
	\sum_{Q \in \Gamma_n} 2^{-C_n(Q)/2} \leq 1,
\end{equation}
which is strictly stronger than the Kraft inequality \eqref{eq:kraft}: by
Cauchy--Schwarz $\sum 2^{-C_n(Q)} \leq \big(\sum 2^{-C_n(Q)/2}\big)^2$, so
\eqref{eq:half_kraft} implies \eqref{eq:kraft} but not conversely. Condition
\eqref{eq:half_kraft} is always achievable at the cost of a factor $2$ in the
penalty: if $\tilde C_n$ satisfies Kraft \eqref{eq:kraft}, then
$C_n := 2\tilde C_n$ satisfies \eqref{eq:half_kraft}, and the oracle inequality
\eqref{eq:main_result} holds with $C_n(Q)/n = 2\tilde C_n(Q)/n$. We state the
theorem under \eqref{eq:half_kraft}; the complexity functions used in
Section~\ref{sec:applications} are defined with this factor absorbed.
\end{remark}

%------------------------------------------------------------------------------
\subsection{Interpretation}
%------------------------------------------------------------------------------

The oracle inequality \eqref{eq:main_result} decomposes the estimation error 
into two terms:

\begin{enumerate}
\item \textbf{Approximation error} $H^2(P, Q)$: How well the model class 
$\Gamma_n$ can approximate the true distribution $P$. This decreases as 
$\Gamma_n$ becomes richer.

\item \textbf{Complexity penalty} $C_n(Q)/n$: The cost of model complexity, 
scaled by sample size. This increases with model richness.
\end{enumerate}

The MDL estimator automatically balances these competing terms, achieving 
a tradeoff that is optimal up to logarithmic factors.

\begin{remark}[On the $1/n$ scaling]
The factor $1/n$ in front of $C_n(Q)$ arises naturally: the complexity 
$C_n(Q)$ represents the total code length for the model, while $n$ observations 
provide $n$ units of information. The ratio measures the per-observation 
cost of model complexity.
\end{remark}

%==============================================================================
\section{Applications to Specific Process Classes}
\label{sec:applications}
%==============================================================================

We now instantiate the general MDL framework for each of the five process 
classes introduced in Section~\ref{sec:intro}.

%------------------------------------------------------------------------------
\subsection{Density Estimation}
\label{subsec:density}
%------------------------------------------------------------------------------

Let $X_1, \ldots, X_n$ be i.i.d.\ observations from an unknown density $p$ 
with respect to a dominating measure $\mu$ on $(E, \mathcal{E})$.

\paragraph{Reference measure.} $P_0 = \mu^n$ (product of dominating measures).

\paragraph{Likelihood.} For candidate density $q \in \Gamma_n$:
\begin{equation}\label{eq:density_likelihood}
\frac{dQ^n}{d\mu^n}(X_1, \ldots, X_n) = \prod_{i=1}^n q(X_i).
\end{equation}

\paragraph{Model class.} Histogram densities on partitions:
\begin{equation}\label{eq:histogram_class}
\Gamma_n^{\text{hist}} = \bigcup_{m=1}^{n} \left\{ q : q = \sum_{j=1}^m 
\frac{\theta_j}{\mu(A_j)} \mathbf{1}_{A_j}, \; \theta_j \in \mathcal{Q}_m, 
\; \sum_j \theta_j = 1 \right\},
\end{equation}
where $\{A_1, \ldots, A_m\}$ is a partition and $\mathcal{Q}_m$ is a finite 
grid of probabilities.

\paragraph{Complexity.}
\begin{equation}\label{eq:density_complexity}
C_n(q) = \log^* m + (m-1) \log m,
\end{equation}
where $\log^* m = \log m + \log \log m + \cdots$ (sum of positive terms).

\paragraph{MDL estimator.}
\begin{equation}\label{eq:density_mdl}
\hat{p}_n = \arg\min_{q \in \Gamma_n^{\text{hist}}} \left( 
-\sum_{i=1}^n \log q(X_i) + C_n(q) \right).
\end{equation}

\paragraph{Hellinger relationship.} By Lemma~\ref{lem:product_measures}:
\begin{equation}\label{eq:density_hellinger}
H^2(P^n, Q^n) = 1 - (1 - H^2(P, Q))^n,
\end{equation}
where $H^2(P, Q) = \frac{1}{2} \int (\sqrt{p} - \sqrt{q})^2 d\mu$.

\begin{theorem}[Density Estimation Convergence]\label{thm:density}
The MDL estimator \eqref{eq:density_mdl} satisfies:
\begin{equation}\label{eq:density_rate}
	H^2(p, \hat{p}_n) \leq \inf_{m \geq 1} \left( \inf_{q \in \Gamma_m} H^2(p, q) 
	+ \frac{\log^* m + (m-1) \log m}{n} \right) + O\left( \frac{\log n}{n} \right)
\end{equation}
$P$-almost surely.
\end{theorem}

%------------------------------------------------------------------------------
\subsection{Poisson Process Estimation}
\label{subsec:poisson}
%------------------------------------------------------------------------------

Let $(T_1, \ldots, T_N)$ be event times from a Poisson process on $E = [0, 1]$ 
with unknown mean measure $\mu$ having intensity $\alpha = d\mu/d\lambda$ 
with respect to Lebesgue measure $\lambda$.

\paragraph{Reference measure.} $P_0$: homogeneous Poisson process with unit 
intensity.

\paragraph{Likelihood.} For candidate intensity $\beta$:
\begin{equation}\label{eq:poisson_likelihood}
\frac{dP_\beta}{dP_0}(T_1, \ldots, T_N) = \exp\left( 
\int_0^1 (1 - \beta(s)) ds \right) \prod_{i=1}^N \beta(T_i).
\end{equation}

\paragraph{Model class.} Piecewise constant intensities:
\begin{equation}\label{eq:poisson_class}
\Gamma_n^{\text{Poisson}} = \bigcup_{m=1}^{n} \left\{ \beta : 
\beta(t) = \sum_{j=1}^m \beta_j \mathbf{1}_{[(j-1)/m, j/m)}(t), \; 
\beta_j \in \mathcal{G}_m \right\},
\end{equation}
where $\mathcal{G}_m = \{k/m : k = 1, \ldots, mn\}$.

\paragraph{Complexity.}
\begin{equation}\label{eq:poisson_complexity}
C_n(\beta) = \log^* m + m \log(mn).
\end{equation}

\paragraph{MDL estimator.}
\begin{equation}\label{eq:poisson_mdl}
\hat{\alpha}_n = \arg\min_{\beta \in \Gamma_n^{\text{Poisson}}} \left( 
\int_0^1 \beta(s) ds - \sum_{i=1}^N \log \beta(T_i) + C_n(\beta) \right).
\end{equation}

\paragraph{Hellinger relationship.}

\begin{lemma}[Hellinger Distance for Poisson Processes]\label{lem:poisson_hellinger}
Let $P_\alpha$ and $P_\beta$ be distributions of Poisson processes with 
intensities $\alpha$ and $\beta$. Then:
\begin{equation}\label{eq:poisson_hellinger}
	H^2(P_\alpha, P_\beta) = 1 - \exp\left( -H^2(\alpha, \beta) \right),
\end{equation}
where $H^2(\alpha, \beta) = \frac{1}{2} \int_0^1 (\sqrt{\alpha(s)} - \sqrt{\beta(s)})^2 ds$.
\end{lemma}

\begin{theorem}[Poisson Process Convergence]\label{thm:poisson}
The MDL estimator \eqref{eq:poisson_mdl} satisfies:
\begin{equation}\label{eq:poisson_rate}
	H^2(\alpha, \hat{\alpha}_n) \leq \inf_{\beta \in \Gamma_n^{\text{Poisson}}} 
	\left( H^2(\alpha, \beta) + \frac{C_n(\beta)}{n} \right) + O\left( \frac{\log n}{n} \right)
\end{equation}
$P$-almost surely, where $n = \mathbb{E}[N] = \int_0^1 \alpha(t) dt$.
\end{theorem}

%------------------------------------------------------------------------------
\subsection{Point Process with Multiplicative Intensity}
\label{subsec:point_process}
%------------------------------------------------------------------------------

Let $N(t)$, $t \in [0, T]$, be a counting process with intensity 
$\lambda(t) = \alpha(t) Y(t)$, where $Y(t)$ is an observable, predictable 
process and $\alpha$ is the unknown baseline intensity \citep{Andersen1993}.

\paragraph{Reference measure.} $P_0$: the distribution when $\alpha \equiv 1$.

\paragraph{Likelihood.} For candidate baseline intensity $\beta$:
\begin{equation}\label{eq:pp_likelihood}
\frac{dP_\beta}{dP_0} = \exp\left( \int_0^T (1 - \beta(s)) Y(s) ds 
+ \int_0^T \log \beta(s) \, dN(s) \right).
\end{equation}

\paragraph{Model class.} Piecewise constant baseline intensities on dyadic 
partitions:
\begin{equation}\label{eq:pp_class}
\Gamma_n^{\text{PP}} = \bigcup_{k=0}^{\lfloor \log_2 n \rfloor} \left\{ 
\beta : \beta(t) = \sum_{j=1}^{2^k} \beta_j \mathbf{1}_{I_{kj}}(t), \; 
\beta_j \in \mathcal{G}_{2^k} \right\},
\end{equation}
where $I_{kj} = [(j-1)T/2^k, jT/2^k)$.

\paragraph{Complexity.}
\begin{equation}\label{eq:pp_complexity}
C_n(\beta) = k + 2^k \log n \quad \text{for } \beta \text{ with } 2^k \text{ pieces}.
\end{equation}

\paragraph{MDL estimator.}
\begin{equation}\label{eq:pp_mdl}
\hat{\alpha}_n = \arg\min_{\beta \in \Gamma_n^{\text{PP}}} \left( 
\int_0^T \beta(s) Y(s) ds - \int_0^T \log \beta(s) \, dN(s) + C_n(\beta) \right).
\end{equation}

\paragraph{Hellinger relationship.}

\begin{lemma}[Hellinger Distance for Point Processes]\label{lem:pp_hellinger}
Let $P_\alpha$ and $P_\beta$ be point process distributions with baseline 
intensities $\alpha$ and $\beta$. Then:
\begin{equation}\label{eq:pp_hellinger}
	H^2(P_\alpha, P_\beta) = 1 - \mathbb{E}_{P_{\sqrt{\alpha\beta}}}\left[ 
	\exp\left( -\frac{1}{2} \int_0^T (\sqrt{\alpha(s)} - \sqrt{\beta(s)})^2 Y(s) ds 
	\right) \right].
\end{equation}
\end{lemma}

\begin{theorem}[Point Process Convergence]\label{thm:pp}
Under the condition $\inf_{t} \mathbb{E}[Y(t)] \geq c > 0$, the MDL estimator 
\eqref{eq:pp_mdl} satisfies:
\begin{equation}\label{eq:pp_rate}
	\int_0^T (\sqrt{\alpha(t)} - \sqrt{\hat{\alpha}_n(t)})^2 dt 
	\lesssim \inf_{\beta \in \Gamma_n^{\text{PP}}} \left( 
	\int_0^T (\sqrt{\alpha} - \sqrt{\beta})^2 dt + \frac{C_n(\beta)}{n} \right)
\end{equation}
$P$-almost surely.
\end{theorem}

%------------------------------------------------------------------------------
\subsection{Markov Transition Model}
\label{subsec:markov}
%------------------------------------------------------------------------------

Consider $n$ independent realizations of a continuous-time Markov chain on 
state space $\{1, \ldots, M\}$ with transition intensities $\alpha^{hj}(t)$ 
for $h \neq j$ \citep{Aalen1978}.

\paragraph{Reference measure.} $P_0$: unit transition intensities.

\paragraph{Likelihood.} For candidate intensities $\gamma = (\gamma^{hj})$:
\begin{equation}\label{eq:markov_likelihood}
\frac{dP_\gamma}{dP_0} = \prod_{h \neq j} \exp\left( 
\int_0^T (1 - \gamma^{hj}(s)) Y_h^{(n)}(s) ds 
+ \int_0^T \log \gamma^{hj}(s) \, dN_{hj}^{(n)}(s) \right),
\end{equation}
where $Y_h^{(n)} = \sum_{i=1}^n Y_h^{(i)}$ and $N_{hj}^{(n)} = \sum_{i=1}^n N_{hj}^{(i)}$.

\paragraph{Model class.} Piecewise constant transition intensities:
\begin{equation}\label{eq:markov_class}
\Gamma_n^{\text{Markov}} = \left\{ \gamma : \gamma^{hj}(t) = 
\sum_{\ell=1}^{m} \gamma^{hj}_\ell \mathbf{1}_{I_\ell}(t), \; 
\gamma^{hj}_\ell \in \mathcal{G}_m \right\}.
\end{equation}

\paragraph{Complexity.}
\begin{equation}\label{eq:markov_complexity}
C_n(\gamma) = \log^* m + M(M-1) \cdot m \cdot \log n.
\end{equation}

\paragraph{MDL estimator.}
\begin{equation}\label{eq:markov_mdl}
\hat{\alpha}_n = \arg\min_{\gamma \in \Gamma_n^{\text{Markov}}} \left( 
\sum_{h \neq j} \left[ \int_0^T \gamma^{hj}(s) Y_h^{(n)}(s) ds 
- \int_0^T \log \gamma^{hj}(s) \, dN_{hj}^{(n)}(s) \right] + C_n(\gamma) \right).
\end{equation}

\paragraph{Hellinger relationship.}

\begin{lemma}[Hellinger Distance for Markov Chains]\label{lem:markov_hellinger}
\begin{equation}\label{eq:markov_hellinger}
	H^2(P_\alpha, P_\gamma) = 1 - \mathbb{E}_{P_{\sqrt{\alpha\gamma}}}\left[ 
	\exp\left( -\frac{1}{2} \sum_{h \neq j} \int_0^T 
	(\sqrt{\alpha^{hj}(s)} - \sqrt{\gamma^{hj}(s)})^2 Y_h(s) ds \right) \right].
\end{equation}
\end{lemma}

%------------------------------------------------------------------------------
\subsection{Diffusion Process Drift Estimation}
\label{subsec:diffusion}
%------------------------------------------------------------------------------

Let $X_t$ satisfy $dX_t = \theta(X_t) dt + \sigma dW_t$ on $[0, T]$, with 
$\sigma > 0$ known \citep{Kutoyants2004}.

\paragraph{Reference measure.} $P_0$: Wiener measure ($\theta \equiv 0$).

\paragraph{Likelihood.} By Girsanov's theorem, for candidate drift $\gamma$:
\begin{equation}\label{eq:diffusion_likelihood}
\frac{dP_\gamma}{dP_0} = \exp\left( \frac{1}{\sigma^2} \int_0^T 
\gamma(X_t) dX_t - \frac{1}{2\sigma^2} \int_0^T \gamma^2(X_t) dt \right).
\end{equation}

\paragraph{Model class.} Piecewise linear drift functions:
\begin{equation}\label{eq:diffusion_class}
\Gamma_n^{\text{Diff}} = \bigcup_{m=1}^{n} \left\{ \gamma : 
\gamma(x) = \sum_{j=1}^{m} (a_j + b_j x) \mathbf{1}_{I_j}(x), \; 
a_j, b_j \in \mathcal{G}_m \right\}.
\end{equation}

\paragraph{Complexity.}
\begin{equation}\label{eq:diffusion_complexity}
C_n(\gamma) = \log^* m + 2m \log n.
\end{equation}

\paragraph{MDL estimator.}
\begin{equation}\label{eq:diffusion_mdl}
\hat{\theta}_n = \arg\min_{\gamma \in \Gamma_n^{\text{Diff}}} \left( 
\frac{1}{2\sigma^2} \int_0^T \gamma^2(X_t) dt 
- \frac{1}{\sigma^2} \int_0^T \gamma(X_t) dX_t + C_n(\gamma) \right).
\end{equation}

\paragraph{Hellinger relationship.}

\begin{lemma}[Hellinger Distance for Diffusions]\label{lem:diffusion_hellinger}
\begin{equation}\label{eq:diffusion_hellinger}
	H^2(P_\theta, P_\gamma) = 1 - \mathbb{E}_{P_{\frac{\theta+\gamma}{2}}}\left[ 
	\exp\left( -\frac{1}{8\sigma^2} \int_0^T (\theta(X_t) - \gamma(X_t))^2 dt 
	\right) \right].
\end{equation}
\end{lemma}

\begin{theorem}[Diffusion Convergence]\label{thm:diffusion}
With $n$ independent replications on $[0, T]$, the MDL estimator 
\eqref{eq:diffusion_mdl} satisfies:
\begin{equation}\label{eq:diffusion_rate}
	\frac{1}{T} \int_0^T (\theta(X_t) - \hat{\theta}_n(X_t))^2 dt 
	\lesssim \inf_{\gamma \in \Gamma_n^{\text{Diff}}} \left( 
	\frac{1}{T} \int_0^T (\theta - \gamma)^2 dt + \frac{C_n(\gamma)}{nT} \right)
\end{equation}
$P$-almost surely.
\end{theorem}

%------------------------------------------------------------------------------
\subsection{Application: Universal Data Compression}
\label{subsec:compression}
%------------------------------------------------------------------------------

The MDL framework provides a principled approach to data compression that 
adapts to the underlying stochastic structure.

\subsubsection{Two-Stage Coding}

A two-stage code for data $\mathbf{X}$ consists of:
\begin{enumerate}
\item \textbf{Model code:} Encode model $\hat{Q}$ using $C_n(\hat{Q})$ bits.
\item \textbf{Data code:} Encode data given model using 
$-\log_2 \hat{Q}(\mathbf{X})$ bits.
\end{enumerate}

Total code length: $L(\mathbf{X}) = -\log_2 \hat{Q}(\mathbf{X}) + C_n(\hat{Q})$.

\begin{theorem}[Universal Compression]\label{thm:compression}
The MDL compressor achieves expected code length:
\begin{equation}\label{eq:compression_bound}
	\mathbb{E}_P[L(\mathbf{X})] \leq H(P) + \inf_{Q \in \Gamma} \left( 
	D(P \| Q) + C(Q) \right) + O(1),
\end{equation}
where $H(P)$ is the entropy of $P$.
\end{theorem}

\subsubsection{Universal Compression Algorithm}

Algorithm~\ref{alg:compression} presents a universal compression scheme 
handling heterogeneous data.

\begin{algorithm}[t]
\caption{Universal MDL Compression}
\label{alg:compression}
\begin{algorithmic}[1]
	\REQUIRE Data sequence $\mathbf{X}$, maximum complexity $K$
	\ENSURE Compressed representation
	
	\STATE $\Gamma \leftarrow \emptyset$
	\STATE \textbf{// Add models from each process class}
	\FOR{each process class $\mathcal{C} \in \{\text{Density, Poisson, PP, Markov, Diffusion}\}$}
	\STATE Add models from $\mathcal{C}$ with complexity up to $K$ to $\Gamma$
	\STATE Adjust complexity: $C(Q) \leftarrow C_{\mathcal{C}}(Q) + \log 5$
	\ENDFOR
	\STATE \textbf{// MDL selection}
	\STATE $\hat{Q} \leftarrow \arg\min_{Q \in \Gamma} (-\log Q(\mathbf{X}) + C(Q))$
	\STATE \textbf{// Encode}
	\STATE ModelCode $\leftarrow$ Encode($\hat{Q}$)
	\STATE DataCode $\leftarrow$ ArithmeticEncode($\mathbf{X}$, $\hat{Q}$)
	\RETURN (ModelCode, DataCode)
\end{algorithmic}
\end{algorithm}

%------------------------------------------------------------------------------
\subsection{Application: Automatic Model Selection}
\label{subsec:model_selection}
%------------------------------------------------------------------------------

When the data-generating mechanism is unknown, MDL provides principled 
model selection.

\begin{algorithm}[t]
\caption{Automatic Model Class Selection}
\label{alg:selection}
\begin{algorithmic}[1]
	\REQUIRE Data $\mathbf{X}$, candidate classes $\mathcal{C}_1, \ldots, \mathcal{C}_K$
	\ENSURE Selected class $\hat{k}$, model $\hat{Q}$
	
	\FOR{$k = 1$ to $K$}
	\STATE $\hat{Q}_k \leftarrow \arg\min_{Q \in \Gamma_n^{(k)}} 
	(-\log Q(\mathbf{X}) + C_n^{(k)}(Q))$
	\STATE $L_k \leftarrow -\log \hat{Q}_k(\mathbf{X}) + C_n^{(k)}(\hat{Q}_k) + \log K$
	\ENDFOR
	\STATE $\hat{k} \leftarrow \arg\min_k L_k$
	\RETURN $(\hat{k}, \hat{Q}_{\hat{k}})$
\end{algorithmic}
\end{algorithm}

\begin{theorem}[Model Selection Consistency]\label{thm:selection}
If $P$ belongs to class $\mathcal{C}_{k^*}$ and is separated from other 
classes ($\inf_{Q \in \Gamma^{(k)}} H^2(P, Q) \geq \delta > 0$ for 
$k \neq k^*$), then Algorithm~\ref{alg:selection} satisfies 
$\mathbb{P}(\hat{k} = k^*) \to 1$.
\end{theorem}

%==============================================================================
\section{Proofs}
\label{sec:proofs}
%==============================================================================

%------------------------------------------------------------------------------
\subsection{Notation and Conventions}
%------------------------------------------------------------------------------

Throughout this section:
\begin{itemize}
\item $H^2(P, Q) = 1 - \rho(P, Q)$ where 
$\rho(P, Q) = \int \sqrt{dP/d\mu \cdot dQ/d\mu} \, d\mu$
\item $C_n: \Gamma_n \to [0, \infty)$ satisfies $\sum_{Q \in \Gamma_n} 2^{-C_n(Q)} \leq 1$
\item $(x)_+ = \max(x, 0)$
\end{itemize}

%------------------------------------------------------------------------------
\subsection{Technical Lemmas}
%------------------------------------------------------------------------------

\begin{lemma}[Likelihood Ratio Bound]\label{lem:likelihood_bound}
For probability measures $P$, $S$ with $P \ll S$, and any $a > 0$:
\begin{equation}\label{eq:Dn_bound}
	P\left( \frac{dP}{dS} \leq a \right) \leq \min\left\{ a, \sqrt{a} \cdot (1 - H^2(P, S)) \right\}.
\end{equation}
\end{lemma}

\begin{proof}
\textbf{First bound:} By Markov's inequality,
\[
P\left( \frac{dS}{dP} \geq \frac{1}{a} \right) \leq a \cdot \mathbb{E}_P\left[ \frac{dS}{dP} \right] = a.
\]

\textbf{Second bound:} Applying Markov to $\sqrt{dS/dP}$:
\[
P\left( \frac{dP}{dS} \leq a \right) \leq \sqrt{a} \cdot \mathbb{E}_P\left[ \sqrt{\frac{dS}{dP}} \right] 
= \sqrt{a} \cdot \rho(P, S) = \sqrt{a} \cdot (1 - H^2(P, S)).
\]
\end{proof}

\begin{lemma}[Penalized Likelihood Comparison]\label{lem:comparison}
Let $Q, S \in \Gamma_n$ and define
\[
A(Q, S) = \left\{ \frac{dQ}{dS} \geq 2^{C_n(Q) - C_n(S)} \right\}.
\]
Then for any $a > 0$:
\begin{equation}\label{eq:comparison_bound}
	P\left( A(Q, S) \cap \left\{ \frac{dP}{dS} \leq a \right\} \right) 
	\leq a \cdot 2^{\frac{1}{2}(C_n(S) - C_n(Q))}.
\end{equation}
\end{lemma}

\begin{proof}
Using exponential bounds with $h = 1/2$:
\begin{align*}
	P\left( A(Q, S) \cap \left\{ \frac{dP}{dS} \leq a \right\} \right)
	&\leq 2^{\frac{1}{2}(C_n(S) - C_n(Q))} \int_{\{dP/dS \leq a\}} \sqrt{\frac{dQ}{dS}} \, dP \\
	&\leq 2^{\frac{1}{2}(C_n(S) - C_n(Q))} \cdot a \int \sqrt{\frac{dQ}{dS}} \, dS \\
	&= a \cdot 2^{\frac{1}{2}(C_n(S) - C_n(Q))}.
\end{align*}
\end{proof}

%------------------------------------------------------------------------------
\subsection{Proof of Main Theorem}
%------------------------------------------------------------------------------

\begin{proof}[Proof of Theorem~\ref{thm:main}]
Let $r_n > 0$, $P_n \in \Gamma_n$, and $a_n = n^{-(1+\epsilon)}$ for $\epsilon > 0$.

Define:
\begin{itemize}
	\item $U_n = \{Q \in \Gamma_n : H^2(P, Q) > r_n\}$
	\item $D_n = \{dP/dP_n \leq a_n\}$
\end{itemize}

\textbf{Step 1: Event decomposition.}
\[
P(\hat{P}_n \in U_n) \leq P\left( \bigcup_{Q \in U_n} A(Q, P_n) \cap D_n \right) + P(D_n^c).
\]

\textbf{Step 2: First term.} By Lemma~\ref{lem:comparison} and union bound:
\begin{align*}
	P\left( \bigcup_{Q \in U_n} A(Q, P_n) \cap D_n \right)
	&\leq \sum_{Q \in U_n} a_n \cdot 2^{\frac{1}{2}(C_n(P_n) - C_n(Q))} \\
	&\leq a_n \cdot 2^{\frac{C_n(P_n)}{2}} \sum_{Q \in \Gamma_n} 2^{-\frac{C_n(Q)}{2}} \\
	&\leq a_n \cdot 2^{\frac{C_n(P_n)}{2}}.
\end{align*}

\textbf{Step 3: Second term.} This step requires structural assumptions on
$(P, P_n)$. Suppose $P = p^{\otimes n}$ and $P_n = p_n^{\otimes n}$ are
$n$-fold product measures (the i.i.d.\ / replicated-sample setting of
Section~\ref{sec:applications}), so that the affinity factorises,
$\rho(P, P_n) = (1 - H^2(p, p_n))^n$ by Lemma~\ref{lem:product_measures}.
Applying the second bound of Lemma~\ref{lem:likelihood_bound} with $S = P_n$
and $a = a_n$ gives
\[
P(D_n) = P\!\left( \frac{dP}{dP_n} \leq a_n \right)
\leq \sqrt{a_n}\,\rho(P, P_n) = \sqrt{a_n}\,(1 - H^2(p, p_n))^n,
\]
and the complementary ``typical-set'' bound used below,
\[
P(D_n^c) \leq \frac{(1 - H^2(p, p_n))^n}{\sqrt{a_n}},
\]
follows by the same Cram\'er--Chernoff control of the log-likelihood ratio
$\log(dP_n/dP)$, whose cumulant generating function is governed by the
per-coordinate affinity. This is the Barron--Cover typical-set argument
\citep{Barron1991}; it uses the product (more generally, factorising-affinity)
structure in an essential way.

\begin{remark}[Scope of Step~3]\label{rem:step3}
The factorisation of the affinity, and hence the exponential bound on
$P(D_n^c)$, holds verbatim only when the data decomposes into $n$ independent
(or independent replicated) components, as in density estimation, the
$n$-replication Markov and diffusion models, and the high-intensity Poisson
regime. For the genuinely continuous-time single-path models (point process
and single-trajectory diffusion) the same conclusion requires a martingale
exponential-inequality substitute for Lemma~\ref{lem:product_measures}; we use
$n$ to denote the relevant information amount ($\mathbb{E}[N]$ or the number of
replications) and treat those cases under the corresponding independence
hypotheses stated in Section~\ref{sec:applications}. The argument should
therefore be read as a Barron--Cover-type sketch under factorising affinity,
not as a single proof covering all five classes uniformly without further
hypotheses.
\end{remark}

\textbf{Step 4: Combining.}
\[
P(\hat{P}_n \in U_n) \leq n^{-(1+\epsilon)} \cdot 2^{\frac{C_n(P_n)}{2}} 
+ n^{\frac{1+\epsilon}{2}} \cdot e^{-n H^2(P, P_n)}.
\]

\textbf{Step 5: Choice of $r_n$.} Set 
$r_n = H^2(P, P_n) + \frac{C_n(P_n)}{n} + \frac{(2+\epsilon)\log n}{n}$.

Both terms are summable, so by Borel-Cantelli:
\[
H^2(P, \hat{P}_n) \leq H^2(P, P_n) + \frac{C_n(P_n)}{n} + O\left(\frac{\log n}{n}\right) 
\quad P\text{-a.s.}
\]

Optimizing over $P_n \in \Gamma_n$ yields the result.
\end{proof}

%------------------------------------------------------------------------------
\subsection{Proofs of Hellinger Relationships}
%------------------------------------------------------------------------------

\begin{proof}[Proof of Lemma~\ref{lem:poisson_hellinger}]
Let $P_0$ be a Poisson process with unit intensity. The Hellinger affinity is:
\begin{align*}
	\rho(P_\alpha, P_\beta) &= \mathbb{E}_{P_0}\left[ \sqrt{\frac{dP_\alpha}{dP_0} 
		\frac{dP_\beta}{dP_0}} \right] \\
	&= \mathbb{E}_{P_0}\left[ e^{1 - \frac{1}{2}(\mu_\alpha + \mu_\beta)} 
	\prod_{i=1}^N \sqrt{\alpha(T_i)\beta(T_i)} \right],
\end{align*}
where $\mu_\alpha = \int \alpha \, ds$.

Using the Laplace functional of Poisson processes:
\begin{align*}
	\rho(P_\alpha, P_\beta) &= \exp\left( \int_0^1 \sqrt{\alpha(s)\beta(s)} \, ds 
	- \frac{1}{2} \int_0^1 (\alpha(s) + \beta(s)) ds \right) \\
	&= \exp\left( -\frac{1}{2} \int_0^1 (\sqrt{\alpha(s)} - \sqrt{\beta(s)})^2 ds \right) \\
	&= \exp(-H^2(\alpha, \beta)).
\end{align*}

Thus $H^2(P_\alpha, P_\beta) = 1 - e^{-H^2(\alpha, \beta)}$.
\end{proof}

\begin{proof}[Proof of Lemma~\ref{lem:pp_hellinger}]
The likelihood is $L_\alpha = \exp(\int (1-\alpha) Y \, ds + \int \log \alpha \, dN)$.

Computing $\sqrt{L_\alpha L_\beta}$:
\begin{align*}
	\sqrt{L_\alpha L_\beta} &= \exp\left( \int (1 - \zeta) Y \, ds + \int \log \zeta \, dN \right) \\
	&\quad \times \exp\left( -\frac{1}{2} \int (\sqrt{\alpha} - \sqrt{\beta})^2 Y \, ds \right),
\end{align*}
where $\zeta = \sqrt{\alpha\beta}$.

The first factor is $L_\zeta$. Taking expectations under $P_0$:
\[
\rho(P_\alpha, P_\beta) = \mathbb{E}_{P_\zeta}\left[ 
\exp\left( -\frac{1}{2} \int (\sqrt{\alpha} - \sqrt{\beta})^2 Y \, ds \right) \right].
\]
\end{proof}

\begin{proof}[Proof of Lemma~\ref{lem:diffusion_hellinger}]
By Girsanov, $L_\theta = \exp(\frac{1}{\sigma^2} \int \theta \, dX - \frac{1}{2\sigma^2} \int \theta^2 \, dt)$.

Let $\zeta = (\theta + \gamma)/2$. Then:
\begin{align*}
	\sqrt{L_\theta L_\gamma} &= L_\zeta \cdot \exp\left( \frac{1}{8\sigma^2} \int (\theta + \gamma)^2 dt 
	- \frac{1}{4\sigma^2} \int (\theta^2 + \gamma^2) dt \right) \\
	&= L_\zeta \cdot \exp\left( -\frac{1}{8\sigma^2} \int (\theta - \gamma)^2 dt \right).
\end{align*}

Thus:
\[
\rho(P_\theta, P_\gamma) = \mathbb{E}_{P_\zeta}\left[ 
\exp\left( -\frac{1}{8\sigma^2} \int (\theta - \gamma)^2 dt \right) \right].
\]
\end{proof}

%==============================================================================
\section{Discussion}
\label{sec:discussion}
%==============================================================================

%------------------------------------------------------------------------------
\subsection{Summary of Contributions}
%------------------------------------------------------------------------------

We have developed a unified framework for statistical estimation across 
multiple stochastic process classes using MDL. The main contributions are:

\begin{enumerate}[label=(\roman*)]
\item A single MDL estimator achieving consistent estimation across 
densities, Poisson processes, point processes, Markov chains, and diffusions.

\item Explicit Hellinger distance formulas sharing a common structure:
\[
H^2(P_\theta, P_\gamma) = 1 - \mathbb{E}_\zeta\left[ 
\exp\left( -c \int d^2(\theta, \gamma) \, d\nu \right) \right].
\]

\item A unified oracle inequality demonstrating automatic bias-complexity 
tradeoff.

\item Practical algorithms for universal compression and model selection.
\end{enumerate}

%------------------------------------------------------------------------------
\subsection{Limitations}
%------------------------------------------------------------------------------

\paragraph{Rates of convergence.}
The rates obtained are not always minimax optimal. For specific model classes,
specialized methods may achieve better rates. Our goal was universality, not
optimality for individual classes. The gap is typically a logarithmic factor: it
is an artefact of the description-length penalty (the per-coefficient quantization
cost $\log n$ in, e.g., $C_n(\beta)=k+2^k\log n$), not of the estimation problem,
and it can be removed by a penalized-projection refinement that pays the model's
metric complexity instead (Remark~\ref{rem:remove-log}).

\paragraph{Computational complexity.}
Exact MDL optimization over large model classes may be computationally 
intensive. Greedy approximations (Algorithm~\ref{alg:compression}) work 
well in practice but lack formal guarantees.

\paragraph{Independence assumptions.}
Several results assume independent observations or replications. Extensions 
to dependent data require additional technical machinery.

\paragraph{Model class specification.}
Results assume the true distribution can be approximated within $\Gamma_n$. 
Model misspecification leads to convergence to the best approximation, which 
may differ from the truth.

%------------------------------------------------------------------------------
\subsection{Extensions}
%------------------------------------------------------------------------------

\paragraph{Dependent observations.}
For $\beta$-mixing sequences, the effective sample size becomes 
$n_{\text{eff}} = n / (1 + 2\sum_k \beta_k)$, and the oracle inequality 
adjusts accordingly.

\paragraph{High-dimensional parameters.}
For sparse models, complexity penalties should reflect sparsity:
$C_n(\theta) = \|\theta\|_0 \log(d/\|\theta\|_0) + \|\theta\|_0 \log n$.

\paragraph{Online estimation.}
Sequential MDL processes observations incrementally:
$\hat{Q}_t = \arg\min_Q (-\sum_{i=1}^t \log Q(X_i | X^{i-1}) + C(Q))$.

\paragraph{Bayesian approaches.}
Instead of model selection, Bayesian MDL computes mixture distributions 
with posterior weights $w(Q | X^n) \propto Q(X^n) 2^{-C_n(Q)}$.

%------------------------------------------------------------------------------
\subsection{Open Problems}
%------------------------------------------------------------------------------

\begin{openproblem}[Minimax optimality]
Under what conditions does MDL achieve minimax optimal rates? Can logarithmic
factors be removed?
\end{openproblem}

\begin{remark}[Removing the logarithmic factor: a partial answer]
\label{rem:remove-log}
For the density/sieve case the answer is now known. The $\log n$ in the piecewise
complexity $C_n(\beta)=k+2^k\log n$ (Eq.~\eqref{eq:pp_complexity}) is a
\emph{coding artefact}: it is the per-piece cost of quantizing $2^k$ continuous
coefficients into a prefix code. An estimator that fits those coefficients
\emph{without} quantizing them---the orthogonal $L^2$-projection (histogram /
piecewise-polynomial) estimator onto the $2^k$-piece model---pays instead the
model's metric complexity $2^k/n$, with \emph{no} logarithm; and the penalized
version selecting $k$ by $\mathrm{pen}(k)\asymp 2^k/n$ attains the exact minimax
rate $n^{-2s/(2s+1)}$ \emph{adaptively}, by the Birg\'e--Massart oracle inequality
\citep{Birge1998,Barron1999}. The mechanism is that the geometric growth of the
dimensions $2^k$ makes the selection weights summable with exponents $\asymp 2^k$,
so no $\log$ re-enters. Thus the description-length MDL is \emph{near-minimax}
(optimal up to $\log$), while its projection counterpart is \emph{exactly}
minimax. What remains open is the general-process version (Poisson, Markov,
diffusion): whether a single contrast-based refinement removes the $\log$
uniformly across the process classes treated here.
\end{remark}

\begin{openproblem}[Model selection across classes]
Can separation conditions in Theorem~\ref{thm:selection} be weakened? What 
happens when multiple classes provide equally good approximations?
\end{openproblem}

\begin{openproblem}[General semimartingales]
Can Hellinger formulas be unified under a general semimartingale framework, 
extending to Lévy processes and jump-diffusions?
\end{openproblem}

\begin{openproblem}[Efficient computation]
For which model classes does polynomial-time exact MDL optimization exist? 
What approximation guarantees can greedy algorithms achieve?
\end{openproblem}

%------------------------------------------------------------------------------
\subsection{Concluding Remarks}
%------------------------------------------------------------------------------

The MDL principle provides a powerful approach to statistical inference that 
naturally balances model fit and complexity. By developing this principle 
across multiple stochastic process classes, we have shown that:

\begin{enumerate}
\item A single estimation procedure achieves consistent inference for 
fundamentally different data-generating mechanisms.

\item The Hellinger distance provides a universal metric with explicit 
formulas linking measure-level and parameter-level distances.

\item The framework has immediate applications to universal compression 
and automatic model selection.
\end{enumerate}

While rates may not always be optimal for specific classes, the universality 
of the approach is valuable for modern applications involving heterogeneous 
data sources.

%==============================================================================
\section*{Acknowledgments}
%==============================================================================

The author thanks the anonymous reviewers for their constructive comments 
that significantly improved this manuscript.

%==============================================================================
\section*{Declaration of Competing Interests}
%==============================================================================

The author declares no competing interests.

%==============================================================================
% References
%==============================================================================

\bibliographystyle{plainnat}

\begin{thebibliography}{99}

\bibitem[Aalen(1978)]{Aalen1978}
Aalen, O.O. (1978).
Nonparametric inference for a family of counting processes.
\textit{Ann.\ Statist.} \textbf{6}, 701--726.

\bibitem[Andersen et al.(1993)]{Andersen1993}
Andersen, P.K., Borgan, {\O}., Gill, R.D. and Keiding, N. (1993).
\textit{Statistical Models Based on Counting Processes}.
Springer, New York.

\bibitem[Barron and Cover(1991)]{Barron1991}
Barron, A.R. and Cover, T.M. (1991).
Minimum complexity density estimation.
\textit{IEEE Trans.\ Inform.\ Theory} \textbf{37}, 1034--1056.

\bibitem[Barron et al.(1998)]{Barron1998}
Barron, A., Rissanen, J. and Yu, B. (1998).
The minimum description length principle in coding and modeling.
\textit{IEEE Trans.\ Inform.\ Theory} \textbf{44}, 2743--2760.

\bibitem[Barron et al.(1999)]{Barron1999}
Barron, A., Birg\'e, L. and Massart, P. (1999).
Risk bounds for model selection via penalization.
\textit{Probab.\ Theory Relat.\ Fields} \textbf{113}, 301--413.

\bibitem[Birg\'e(1983)]{Birge1983}
Birg\'e, L. (1983).
Approximation dans les espaces m\'etriques et th\'eorie de l'estimation.
\textit{Z.\ Wahrsch.\ Verw.\ Gebiete} \textbf{65}, 181--237.

\bibitem[Birg\'e and Massart(1998)]{Birge1998}
Birg\'e, L. and Massart, P. (1998).
Minimum contrast estimators on sieves: Exponential bounds and rates of convergence.
\textit{Bernoulli} \textbf{4}, 329--375.

\bibitem[Birg\'e(2006)]{Birge2006}
Birg\'e, L. (2006).
Model selection via testing: An alternative to (penalized) maximum likelihood estimators.
\textit{Ann.\ Inst.\ H.\ Poincar\'e Probab.\ Statist.} \textbf{42}, 273--325.

\bibitem[Cover and Thomas(2006)]{Cover2006}
Cover, T.M. and Thomas, J.A. (2006).
\textit{Elements of Information Theory}, 2nd ed.
Wiley, Hoboken.

\bibitem[Devroye and Lugosi(2001)]{Devroye2001}
Devroye, L. and Lugosi, G. (2001).
\textit{Combinatorial Methods in Density Estimation}.
Springer, New York.

\bibitem[Gr\"unwald(2007)]{Grunwald2007}
Gr\"unwald, P.D. (2007).
\textit{The Minimum Description Length Principle}.
MIT Press, Cambridge.

\bibitem[Hansen and Yu(2001)]{Hansen2001}
Hansen, M.H. and Yu, B. (2001).
Model selection and the principle of minimum description length.
\textit{J.\ Amer.\ Statist.\ Assoc.} \textbf{96}, 746--774.

\bibitem[Kakutani(1948)]{Kakutani1948}
Kakutani, S. (1948).
On equivalence of infinite product measures.
\textit{Ann.\ Math.} \textbf{49}, 214--224.

\bibitem[Karr(1987)]{Karr1987}
Karr, A.F. (1987).
Maximum likelihood estimation in the multiplicative intensity model via sieves.
\textit{Ann.\ Statist.} \textbf{15}, 473--490.

\bibitem[Kutoyants(2004)]{Kutoyants2004}
Kutoyants, Yu.A. (2004).
\textit{Statistical Inference for Ergodic Diffusion Processes}.
Springer, London.

\bibitem[Le Cam(1973)]{LeCam1973}
Le Cam, L. (1973).
Convergence of estimates under dimensionality restrictions.
\textit{Ann.\ Statist.} \textbf{1}, 38--53.

\bibitem[Le Cam(1986)]{LeCam1986}
Le Cam, L. (1986).
\textit{Asymptotic Methods in Statistical Decision Theory}.
Springer, New York.

\bibitem[McKeague(1986)]{McKeague1986}
McKeague, I.W. (1986).
Estimation for a semimartingale regression model using the method of sieves.
\textit{Ann.\ Statist.} \textbf{14}, 579--589.

\bibitem[Prakasa Rao(1999)]{PrakasaRao1999a}
Prakasa Rao, B.L.S. (1999).
\textit{Statistical Inference for Diffusion Type Processes}.
Arnold, London.

\bibitem[Rissanen(1978)]{Rissanen1978}
Rissanen, J. (1978).
Modeling by shortest data description.
\textit{Automatica} \textbf{14}, 465--471.

\bibitem[Rissanen(1983)]{Rissanen1983}
Rissanen, J. (1983).
A universal prior for integers and estimation by minimum description length.
\textit{Ann.\ Statist.} \textbf{11}, 416--431.

\bibitem[Rissanen(1984)]{Rissanen1984}
Rissanen, J. (1984).
Universal coding, information, prediction \& estimation.
\textit{IEEE Trans.\ Inform.\ Theory} \textbf{30}, 629--636.

\bibitem[Tsybakov(2009)]{Tsybakov2009}
Tsybakov, A.B. (2009).
\textit{Introduction to Nonparametric Estimation}.
Springer, New York.

\bibitem[Yang and Barron(1999)]{Yang1999}
Yang, Y. and Barron, A. (1999).
Information-theoretic determination of minimax rates of convergence.
\textit{Ann.\ Statist.} \textbf{27}, 1564--1599.

\end{thebibliography}

\end{document}